\section{Fundamentals}

\subsection{Probability Space}
\begin{definition}[\(\sigma\)-Algebra]
    \label{def:sigma-algebra}%
    Suppose \(\Omega\) is a set, and \(\mathcal{F} \subseteq \powerset\para{\Omega}\) is a collection of subsets of \(\Omega\). We say \(\mathcal{F}\) is a \emph{\(\sigma\)-algebra}, if the following hold:
    \begin{enumerate}
        \item \(\Omega \in \mathcal{F}\);
        \item if \(A \in \mathcal{F}\), then the complement \(A\compl \in \mathcal{F}\);
        \item if \(\para{A_n}_{n \in \NN}\) is a countable collection of sets in \(\mathcal{F}\) (i.e. \(A_n \in \mathcal{F}\) for all \(n \in \NN\)), then
        \[
            \bigcup_{n \in \NN} A_n \in \mathcal{F}.
        \]
    \end{enumerate}
\end{definition}

\begin{definition}[Probability Measure]
    \label{def:probability-measure}%
    Suppose \(\mathcal{F}\) is a \(\sigma\)-algebra on \(\Omega\). A function
    \[
        \functiondomain{\Prob}{\mathcal{F}}{\ccintv{0}{1}}
    \]
    is called a \emph{probability measure} if
    \begin{enumerate}
        \item \(\Prob\para{\Omega} = 1\);
        \item for any countable disjoint collection \(\para{A_n}_{n \in \NN}\) in \(\mathcal{F}\), we have
        \[
            \Prob\para{\bigcup_{n \in \NN} A_n} = \sum_{n \in \NN} \Prob\para{A_n}.
        \]
    \end{enumerate}

    The value of \(\Prob\para{A}\) for some \(A \in \mathcal{F}\) is the \emph{probability of \(A\)}.
\end{definition}

\begin{definition}[Probability Space]
    \label{def:probability-space}%
    For a set \(\Omega\), a \(\sigma\)-algebra \(\mathcal{F}\) on \(\Omega\), and a probability measure \(\Prob\) on \(\mathcal{F}\) and \(\Omega\), the triple
    \[
        \para{\Omega, \mathcal{F}, \Prob}
    \]
    is a \emph{probability space}.
\end{definition}

\begin{definition}[Outcomes, Events]
    \label{def:outcome-event}%
    The elements of the set \(\Omega\) are called \emph{outcomes}, and the elements of the \(\sigma\)-algebra \(\mathcal{F}\) are called \emph{events}.
\end{definition}

\begin{remark}
    By convention, if \(\Omega\) is countable, then we take \(\mathcal{F}\) to be \(\powerset\para{\Omega}\), the set of all subsets of \(\Omega\).
\end{remark}

We note that we talk about probability of \emph{events}, rather than \emph{outcomes} -- the probability function is defined from \(\mathcal{F}\), the \(\sigma\)-algebra.

We may conclude some properties of the probability measure \(\Prob\) from the definition.

\begin{proposition}[Properties of \(\Prob\)]
    \label{prop:properties-prob}%
    \begin{itemize}
        \item \(\Prob\para{A\compl} = 1 - \Prob\para{A}\).
        \item \(\Prob\para{\emptyset} = 0\).
        \item \(\Prob\para{A} \leq \Prob\para{B}\) if \(A \subseteq B\).
        \item \(\Prob\para{A \cup B} = \Prob\para{A} + \Prob\para{B} - \Prob\para{A \cap B}\).
    \end{itemize}
\end{proposition}

\subsection{Equally Likely Outcomes}
Now we consider some examples of probability spaces, where the outcomes are finite, and all outcomes are equally likely.

\begin{example}[Rolling a Fair Die]
    \label{ex:rolling-a-fair-die}%
    We consider the outcomes \(\Omega = \brce{1, 2, \ldots, 6}\), and \(\mathcal{F} = \powerset\para{\Omega}\) for all subsets of \(\Omega\) (since \(\Omega\) is countable).
        
    We naturally take for all \(\omega \in \Omega\),
    \[
        \Prob\para{\brce{\omega}} = \frac{1}{6},
    \]
    and for all \(A \subseteq \Omega\), we take
    \[
        \Prob\para{A} = \frac{\abs{A}}{6}.
    \]

    All outcomes are equally likely.
\end{example}

Lots of situations are like this in real life, and we refer to such setup as \emph{classical probability}, with \emph{equally likely outcomes}.

\begin{definition}[Equally Likely Outcomes]
    \label{def:equally-likely-outcomes}%
    For any finite set \(\Omega = \brce{\omega_1, \omega_2, \ldots, \omega_n}\), take \(\mathcal{F} = \powerset\para{\Omega}\). We define the probability measure to be
    \[
        \function{\Prob}{\mathcal{F}}{\ccintv{0}{1}}{A}{\frac{\abs{A}}{\abs{\Omega}}.}
    \]

    This is called the \emph{uniform probability distribution}.
\end{definition}

This models the experiment of picking a \emph{random} element of \(\Omega\). For all \(\omega \in \Omega\),
\[
    \Prob\para{\brce{\omega}} = \frac{1}{\abs{\Omega}}
\]

We look at some more examples on calculating the probability under such setup. Effectively, those are combinatorial problems counting the size of sets.

This example is more about counting the size of the \emph{sample space} which is the denominator.

\begin{example}[Picking Balls from a Bag]
    \label{ex:picking-balls-from-bag}%
    Suppose we have \(n\) balls labelled \(\brce{1, \ldots, n}\), and they are indistinguishable by touch. Pick \(k \leq n\) balls \emph{at random} (any outcome is equally likely) without replacement. We let the outcomes be
    \[
        \Omega = \set{A \subseteq \brce{1, 2, \ldots, n}}{\abs{A} = k}
    \]
    and hence
    \[
        \abs{\Omega} = \binom{n}{k}.
    \]

    Hence, for any \(\omega \in \Omega\), we have
    \[
        \Prob\para{\brce{\omega}} = \frac{1}{\binom{n}{k}}.
    \]
\end{example}

The following examples also require counting on the numerator, which is the cardinality of the event.

\begin{example}[Deck of Cards]
    \label{ex:deck-of-cards}%
    In a well-shuffled deck of 52 cards, take
    \[
        \Omega = \brce{\text{all permutations of 52 cards}}
    \]
    and hence
    \[
        \abs{\Omega} = 52! = 80,658,175,170,943,878,571,660,636,856,403,766,975,289,505,440,883,277,824,000,000,000,000.
    \]
    
    Take \(\mathcal{F} = \powerset\para{\Omega}\). Then
    \[
        \Prob\para{\text{top 2 cards are aces}} = \frac{4 \times 3 \times 50!}{52!}.
    \]
\end{example}

The next example uses results on set cardinalities.

\begin{example}[Largest Digit]
    \label{ex:largest-digit}%
    Consider a string of \(n\) random digits from \(0, 1, \ldots, 9\). Let
    \[
        \Omega = \brce{0, 1, \ldots, 9}^n,
    \]
    and we have \(\abs{\Omega} = 10^n\).

    Let the event \(A_k = \brce{\text{no digit exceeds }k}\), then
    \[
        \Prob\para{A_k} = \frac{\para{k + 1}^{n}}{10^n}.
    \]

    Now we consider the event \(B_k = \brce{\text{the largest digit is equal to }k}\), then note \(B_k = A_k \setminus A_{k - 1}\), so
    \[
        \Prob\para{B_k} = \Prob\para{A_k \setminus A_{k - 1}} = \Prob\para{A_k} - \Prob\para{A_{k - 1}} = \frac{\para{k + 1}^n - k^n}{10^n}.
    \]
\end{example}

\begin{example}[Birthday Problem]
    \label{ex:birthbay-problem}%
    There are \(n\) people. What is the probability that at least 2 of them share the same birthday?
        
    Assume each birthday is equally likely to be one of \(\brce{1, 2, \ldots, 365}\). Take
    \[
        \Omega = \brce{1, 2, \ldots, 365}^n,
    \]
    and take \(\mathcal{F} = \powerset\para{\Omega}\). Hence, for all \(\omega \in \Omega\), we have
    \[
        \Prob\para{\brce{\omega}} = \frac{1}{365^n}.
    \]

    Consider the event \(A = \brce{\text{at least 2 people have the same birthday}}\).

    Then \(A\compl = \brce{\text{all }n\text{ birthdays are different}}\). Hence,
    \[
        \Prob\para{A\compl} = \frac{365 \times 364 \times \cdots \times \para{365 - n + 1}}{365^n}
    \]
    and hence
    \[
        \Prob\para{A} = 1 - \Prob\para{A\compl}.
    \]

    Note that \(n = 22\) gives \(\Prob\para{A} \approx 0.476\), and \(n = 23\) gives \(\Prob\para{A} \approx 0.507\).
\end{example}

\subsection{Combinatorial Analysis}
We start off by investigating some easy problems in combinatorics.

\begin{example}[Partitioning Sets into Fixed Sizes]
    \label{ex:partition-sets-given-size}%
    We first consider \(\Omega\) being a finite set with \(\abs{\Omega} = n\). We want to partition \(\Omega\) into \(k\) disjoint subsets \(\Omega_1, \Omega_2, ldots, \Omega_k\) with \(\Omega_i = n_i\) for \(\sum_{i = 1}^{k} n_i = n\).

    Let \(M\) be the number of ways to do this. We have
    \begin{align*}
        M &= \binom{n}{n_1} \times \binom{n - n_1}{n_2} \times \cdots \times \binom{n - \para{n_1 + \cdots + n_{k - 1}}}{n_k} \\
        &= \frac{n!}{n_1! n_2! \cdots n_k!}.
    \end{align*}
\end{example}

\begin{definition}[Multinomial Coefficient]
    \label{def:multinomial-coeff}%
    The \emph{multinomial coefficient} is defined as
    \[
        \binom{n}{n_1, n_2, \ldots, n_k} = \frac{n!}{n_1! n_2! \cdots n_k!}.
    \]
\end{definition}

\begin{examples}[Increasing and Strictly Increasing Functions]
    \label{ex:increasing-functions}%
    We say \(f\) is \emph{strictly increasing} if \(x < y\) implies \(f\para{x} < f\para{y}\), and \(f\) is \emph{(weakly) increasing}, if \(x \leq y\) implies \(f\para{x} \leq f\para{y}\).

    Let \(\functiondomain{f}{\brce{1, 2, \ldots, k}}{\brce{1, 2, \ldots, n}}\).
    
    \begin{enumerate}
        \item \textit{Strictly Increasing Functions.} Let \(k \leq n\). How many strictly increasing functions \(f\) are there? Such a function is uniquely determined by its range, so there are \(\binom{n}{k}\) of them.
        \item \textit{Increasing Functions.} How many increasing functions \(f\) are there?
        
        Let the function space
        \[
            \mathcal{F}\para{k, n} = \set{f}{\functiondomain{f}{\brce{1, 2, \ldots, k}}{\brce{1, 2, \ldots, n}}}.
        \]
        
        We define a bijection
        \[
            \set{f \in \mathcal{F}\para{k, n}}{f \text{ is increasing}} \to \set{g \in \mathcal{F}\para{k, n + k - 1}}{g \text{ is strictly increasing}},
        \]
        where \(g\para{i} = f\para{i} + i - 1\) that is strictly increasing. Hence, there are
        \[
            \binom{n + k - 1}{k}
        \]
        such increasing functions in \(\mathcal{F}\para{k, n}\).
    \end{enumerate}
\end{examples}

\subsection{Stirling's Formula}

The following asymptotic formula is very useful in the study of probability. Although the proof is non-examinable, we would still give a proof of it.

\begin{definition}[Asymptotic]
    \label{def:asymptotic}%
    For two sequences \(\para{a_n}, \para{b_n}\), we say \(\para{a_n}\) is \emph{asymptotic} to \(\para{b_n}\), denoted as
    \[
        a_n \sim b_n
    \]
    as \(n \to \infty\), if \(\frac{a_n}{b_n} \to 1\) as \(n \to \infty\).
\end{definition}

\begin{theorem}[Stirling's Formula]
    \label{thm:stirling-formula}%
    We have
    \[
        n! \sim n^n \sqrt{2 \pi n} \cdot e^{-n},
    \]
    as \(n \to \infty\).
\end{theorem}

This statement is relatively long to prove (and is non-examinable), but a statement that is also useful is given below, which is a corollary of the above statement, but whose proof is \textbf{examinable}.

\begin{corollary}
    \label{cor:weaker-statement-of-stirling}%
    We have \(\log \para{n!} \sim n \log n\) as \(n \to \infty\).
\end{corollary}

\begin{proof}
    Define
    \[
        l_n = \log\para{n!} = \log 2 + \cdots + \log n.
    \]

    For all \(x \in \RR\), denote \(\floor{x}\) as the integer part of \(x\). We have
    \[
        \log \floor{x} \leq \log x \leq \log \floor{x + 1}.
    \]

    We integrate this inequality from \(1\) to \(n\), and hence, we have
    \[
        \sum_{k = 1}^{n - 1} \log k \leq \int_{1}^{n} \log x \Diff x \leq \sum_{k = 1}^{n} \log k
    \]
    and hence we have
    \[
        l_{n - 1} \leq n \log n - n + 1 \leq l_n.
    \]

    Therefore, rearranging this gives
    \[
        n \log n - n + 1 \leq l_n \leq \para{n + 1} \log \para{n + 1} - \para{n + 1} + 1,
    \]
    and dividing this through by \(n \log n\) gives
    \[
        \frac{l_n}{n \log n} \to 1
    \]
    as \(n \to \infty\).
\end{proof}

We first show that Stirling's Formula is true, up to some constant (which would turn out to be \(\sqrt{2\pi}\)). A lemma that could be useful is a result that comes from integration by parts.

\begin{lemma}
    \label{lem:some-integration-formula}%
    For any \(\functiondomain{f}{\RR}{\RR}\) which is twice differentiable, for all \(a < b\), we have
    \[
        \int_a^b f\para{x} \Diff x = \frac{f\para{a} + f\para{b}}{2} \cdot \para{b - a} - \frac{1}{2} \int_a^b \para{x - a} \para{b - x} f''\para{x} \Diff x.
    \]
\end{lemma}

\begin{proof}
    Using integration by parts twice.
\end{proof}


\begin{lemma}[Stirling's Formula with Constant]
    \label{lem:stirling-formula-constant}%
    We have
    \[
        n! \sim A n^n \sqrt{n} \cdot e^{-n},
    \]
    as \(n \to \infty\), for some real number \(A \in \RR\).
\end{lemma}

\begin{proof}
    Take \(f\para{x} = \log x\), with \(a = k\), \(b = k + 1\), for \(k \in \NN\). Using \autoref{lem:some-integration-formula}, we have
    \[
        \int_{k}^{k + 1} \log x \Diff x = \frac{\log k + \log\para{k + 1}}{2} + \frac{1}{2} \int_{k}^{k + 1} \frac{\para{x - k} \para{k + 1 - x}}{x^2} \Diff x.
    \]

    Summing this from \(k = 1\) to \(n - 1\) gives
    \[
        \int_{1}^{n} \log x \Diff x = \frac{\log \para{n - 1}! + \log \para{n!}}{2} + \sum_{k = 1}^{n - 1} a_k,
    \]
    where
    \[
        a_k = \frac{1}{2} \int_{0}^{1} \frac{x \para{1 - x}}{\para{x + k}^2} \Diff x.
    \]

    Hence, we have
    \[
        n \log n - n + 1 = \log\para{n!} - \frac{\log n}{2} + \sum_{k = 1}^{n - 1} a_k
    \]

    Therefore,
    \[
        \log \para{n!} = n \log n - n + 1 + \frac{\log n}{2} - \sum_{k = 1}^{n - 1} a_k,
    \]
    and hence
    \[
        n! = n^n \cdot e^{-n} \cdot \sqrt{n} \cdot \exp\para{1 - \sum_{k = 1}^{n - 1} a_k}.
    \]

    Now, we have
    \[
        a_k = \frac{1}{2} \int_{0}^{1} \frac{x \para{1 - x}}{\para{x + k}^2} \Diff x \leq \frac{1}{2k^2} \int_{0}^{1} x\para{1 - x} \Diff x = \frac{1}{12k^2}.
    \]

    Therefore,
    \[
        \sum_{k = 1}^{\infty} a_k < \infty.
    \]

    Let \(A = \exp\para{1 - \sum_{k = 1}^{\infty} a_k}\). Therefore,
    \[
        n! = n^n \cdot e^{-n} \cdot \sqrt{n} \cdot A \cdot \exp\para{\sum_{k = n}^{\infty} a_k}.
    \]

    The \(\exp\para{\sum_{k = n}^{\infty} a_k}\) will converge to \(1\) as \(n \to \infty\), since the inner sum converges to \(0\).

    We have showed that
    \[
        \frac{n!}{n^n \cdot e^{-n} \cdot \sqrt{n}} \to A
    \]
    as \(n \to \infty\), as desired.
\end{proof}

To finish the proof of \autoref{thm:stirling-formula} (Stirling's Formula), we want to find \(A\). An asymptotic analysis on the binomial coefficient \(2^{-2n} \binom{2n}{n}\) would be useful, and one method to do this is to use the expansion of the binomial coefficient.

\begin{lemma}[Asymptotic for \(2^{-2n} \binom{2n}{n}\), using Factorials]
    \label{lem:asymptotic-binomial-factorial}%
    We have
    \[
        2^{-2n} \binom{2n}{n} \sim \frac{\sqrt{2}}{A \sqrt{n}}
    \]
    for the \(A\) above, in \autoref{lem:stirling-formula-constant}.
\end{lemma}

\begin{proof}
    Notice that
    \begin{align*}
        2^{-2n} \binom{2n}{n} &= 2^{-2n} \cdot \frac{\para{2n}!}{n! \cdot n!} \\
        &\sim \frac{2^{-2n} \cdot \para{2n}^{2n} \cdot e^{-2n} \cdot \sqrt{2n} \cdot A}{n^n \cdot e^{-n} \cdot \sqrt{n} \cdot A \cdot n^n \cdot e^{-n} \cdot \sqrt{n} \cdot A} \\
        &= \frac{\sqrt{2}}{A \sqrt{n}}.
    \end{align*}
\end{proof}

Another way to study the asymptotic behaviour of this function, is to relate this with the Wallis Integral.

\begin{lemma}[The Wallis Integrals]
    \label{lem:wallis-integral}%
    Let
    \[
        I_n = \int_{0}^{\frac{\pi}{2}} \cos^n \theta \Diff \theta
    \]
    be the Wallis Integrals. We have
    \[
        I_{2n} = \frac{\pi}{2} \cdot 2^{-2n} \cdot \binom{2n}{n},
    \]
    and
    \[
        I_{2n + 1} = \frac{1}{2n + 1} \cdot \para{2^{-2n} \cdot \binom{2n}{n}}\inverse.
    \]
\end{lemma}

\begin{proof}
    We notice \(I_0 = \frac{\pi}{2}\) and \(I_2 = 1\). Integration by parts gives us
    \[
        I_n = \frac{n - 1}{n} \cdot I_{n - 2},
    \]
    so
    \begin{align*}
        I_{2n} &= \frac{2n - 1}{2n} \cdot I_{2n - 2} \\
        &= \cdots \\
        &= \frac{\para{2n - 1} \cdot \para{2n - 3} \cdots 3 \cdot 1}{\para{2n} \cdot \para{2n - 2} \cdots 2} I_0 \\
        &= \frac{\para{2n} \cdot \para{2n - 1} \cdot \para{2n - 2} \cdot \para{2n - 3} \cdots 3 \cdot 2 \cdot 1}{2n \cdot 2n \cdot \para{2n - 2} \cdot \para{2n - 2} \cdots 2 \cdot 2} \cdot \frac{\pi}{2} \\
        &= \frac{\para{2n}!}{2^{2n} \cdot n! \cdot n!} \cdot \frac{\pi}{2} \\
        &= \frac{\pi}{2} \cdot 2^{-2n} \cdot \binom{2n}{n},
    \end{align*}
    so we have
    \[
        I_{2n} = \frac{\pi}{2} \cdot 2^{-2n} \cdot \binom{2n}{n}.
    \]

    Similarly, we have
    \[
        I_{2n + 1} = \frac{2n \cdot \para{2n - 1} \cdots 4 \cdot 2}{\para{2n + 1} \cdot \para{2n - 1} \cdots 3 \cdot 1} I_1 = \frac{1}{2n + 1} \cdot \para{2^{-2n} \cdot \binom{2n}{n}}\inverse.
    \]

\end{proof}

\begin{corollary}[Asymptotic for \(2^{-2n} \binom{2n}{n}\), using Wallis]
    \label{cor:asymptotic-binomial-wallis}%
    We have
    \[
        2^{-2n} \binom{2n}{n} \sim \frac{1}{\sqrt{\pi n}}.
    \]
\end{corollary}

\begin{proof}
    We use \autoref{lem:wallis-integral}. 

    Notice that we have \(\frac{I_n}{I_{n - 2}} \to 1\) as \(n \to \infty\).

    Note that \(I_n\) is decreasing in \(n\), so
    \[
        \frac{I_{2n}}{I_{2n + 1}} \leq \frac{I_{2n - 1}}{I_{2n + 1}} \to 1
    \]
    and
    \[
        \frac{I_{2n}}{I_{2n + 1}} \geq \frac{I_{2n}}{I_{2n - 1}} \to 1,
    \]
    so we have \(\frac{I_{2n}}{I_{2n + 1}} \to 1\) as \(n \to \infty\).

    Hence,
    \[
        \frac{\frac{\pi}{2} \cdot 2^{-2n} \cdot \binom{2n}{n}}{\frac{1}{2n + 1} \cdot \para{2^{-2n} \cdot \binom{2n}{n}}\inverse} \to 1
    \]
    which gives
    \[
        \para{2^{-2n} \cdot \binom{2n}{n}}^2 \sim \frac{1}{\pi n}
    \]
    as \(n \to \infty\), as desired.
\end{proof}

Now we are ready to find \(A\), and give a proof of the original statement in \autoref{thm:stirling-formula} (Stirling's formula).

\begin{proof}[of \autoref{thm:stirling-formula} (Stirling's Formula)]
    From \autoref{lem:stirling-formula-constant}, it remains to show that \(A = \sqrt{2 \pi}\).
    
    Using the result in \autoref{lem:asymptotic-binomial-factorial}, we have
    \[
        2^{-2n} \binom{2n}{n} \sim \frac{\sqrt{2}}{A \sqrt{n}}.
    \]

    From \autoref{cor:asymptotic-binomial-wallis}, we also have
    \[
        2^{-2n} \binom{2n}{n} \sim \frac{1}{\sqrt{\pi n}}.
    \]
    
    This then implies \(A = \sqrt{2 \pi}\), which shows the result as desired.
\end{proof}

\subsection{Properties of the Probability Measure}

Let \(\para{\Omega, \mathcal{F}, \Prob}\) be a probability space. We give some useful properties of the probability measure \(\Prob\), other than those on the example sheets.

\begin{proposition}[Countable Subadditivity]
    \label{prop:countable-subadditivity}%
    If \(\para{A_n}\) is a collection in \(\mathcal{F}\), then
    \[
        \Prob\para{\bigcup_{n \in \NN} A_n} \leq \sum_{n \in \NN} \Prob\para{A_n}.
    \]
\end{proposition}

\begin{proof}
    Define \(B_1 = A_1\), and for \(n \geq 2\), we let
    \[
        B_n = A_n \setminus \para{A_1 \cup A_2 \cup \cdots \cup A_{n - 1}}.
    \]

    Then \(\para{B_n}\) is a disjoint collection of sets in \(\mathcal{F}\), and we have
    \[
        \bigcup_{n \in \NN} B_n = \bigcup_{n \in \NN} A_n,
    \]
    so by the countable additivity property (see \autoref{def:probability-measure}, second point), we have
    \[
        \Prob\para{\bigcup A_n} = \Prob\para{\bigcup B_n} = \sum \Prob\para{B_n} \leq \sum \Prob\para{A_n}
    \]
    since \(B_n \subseteq A_n\).
\end{proof}

\begin{proposition}[Continuity of Probability Measure]
    \label{prop:continuity-of-probability-measure}%
    For some \emph{increasing} sequence of events \(\para{A_n}_{n \in \NN}\) in \(\mathcal{F}\) (i.e., we have \(A_n \subseteq A_{n + 1}\) for all \(n \in \NN\)), then we have \(\para{\Prob\para{A_n}}\) is increasing in \(n\), and converges to the limit
    \[
        \Prob\para{A_n} \to \Prob\para{\bigcup_{n = 1}^{\infty} A_n}.
    \]
\end{proposition}

\begin{proof}
    We define \(B_1 = A_1\), and for \(n \geq 2\), take
    \[
        B_n = A_n \setminus \bigcup_{m = 1}^{n - 1} A_m.
    \]

    Then \(\para{B_n}\) are disjoint, and we have
    \[
        \bigcup_{k = 1}^{n} B_k = A_n,
    \]
    so
    \[
        \Prob\para{A_n} = \Prob\para{\bigcup_{k = 1}^{n} B_k} = \sum_{k = 1}^{n} \Prob\para{B_k} \to \sum_{k = 1}^{\infty} \Prob\para{B_k}.
    \]

    By the countable additivity property, we have
    \[
        \Prob\para{\bigcup_{k = 1}^{\infty} B_k} = \sum_{k = 1}^{\infty} \Prob\para{B_k},
    \]
    and we also have
    \[
        \bigcup_{k = 1}^{\infty} B_k = \bigcup_{k = 1}^{\infty} A_k,
    \]
    Therefore
    \[
        \sum_{k = 1}^{\infty} \Prob\para{B_k} = \Prob\para{\bigcup_{k = 1}^{\infty} A_k},
    \]
    which finishes our proof.
\end{proof}

\begin{corollary}[Countinuity of Probability Measure for Decreasing Events]
    \label{cor:continuity-probability-measure-decreasing-sequence}%
     For some \emph{decreasing} sequence of events \(\para{A_n}_{n \in \NN}\) in \(\mathcal{F}\) (i.e., we have \(A_n \supseteq A_{n + 1}\) for all \(n \in \NN\)), then we have \(\para{\Prob\para{A_n}}\) is decreasing in \(n\), and converges to the limit
    \[
        \Prob\para{A_n} \to \Prob\para{\bigcap_{n = 1}^{\infty} A_n}.
    \]
\end{corollary}

\begin{proof}
    We apply \autoref{prop:continuity-of-probability-measure} by taking the complements.
\end{proof}

\begin{remark}
    There is \emph{no} topology involved in play here -- this continuity is different, and the definition we gave above is what is meant for this probability measure.
\end{remark}

We recall that in terms of the cardinality (the cardinality measure on finite sets) that we have the Inclusion--Exclusion Principle. A similar formula holds for two events, \(A, B \in \mathcal{F}\), which says
\[
    \Prob\para{A \cup B} = \Prob\para{A} + \Prob\para{B} - \Prob\para{A \cap B}.
\]

Applying this again to three events \(A, B, C \in \mathcal{F}\), we have
\begin{align*}
    \Prob\para{A \cup B \cup C} &=  \Prob\para{\para{A \cup B} \cup C}\\
    &= \Prob\para{A \cup B} + \Prob\para{C} - \Prob\para{\para{A \cap C} \cup \para{B \cap C}} \\
    &= \Prob\para{A} + \Prob\para{B} + \Prob\para{C} - \Prob\para{A \cap B} - \Prob\para{A \cap C} - \Prob\para{B \cap C} + \Prob\para{A \cap B \cap C}.
\end{align*}

This leads to \autoref{thm:inclusion-exclusion}, the Inclusion--Exclusion Formula.

\begin{theorem}[Inclusion--Exclusion Formula]
    \label{thm:inclusion-exclusion}%%
    If we have events \(A_1, A_2, \ldots, \mathcal{F}\), then we have
    \[
        \Prob\para{\bigcup_{i = 1}^{n} A_i} = \sum_{k = 1}^{n} \para{-1}^{k + 1} \sum_{1 \leq i_1 < \cdots < i_k \leq n} \Prob\para{A_{i_1} \cap A_{i_2} \cap \cdots \cap A_{i_k}}.
    \]
\end{theorem}

\begin{proof}
    This holds for \(n = 2\) events. Assume it holds for \(n - 1\) events.

    We have
    \begin{align*}
        \Prob\para{\bigcup_{i = 1}^{n} A_i} &= \Prob\para{\bigcup_{i = 1}^{n - 1} A_i \cup A_n} \\
        &= \Prob\para{\bigcup_{i = 1}^{n - 1} A_i} + \Prob\para{A_n} - \Prob\para{\bigcup_{i = 1}^{n - 1} A_i \cap A_n} \\
        &= \Prob\para{\bigcup_{i = 1}^{n - 1} A_i} + \Prob\para{A_n} - \Prob\para{\bigcup_{i = 1}^{n - 1} \para{A_i \cap A_n}}.
    \end{align*}

    No, we apply the induction hypothesis to the terms
    \[
        \Prob\para{\bigcup_{i = 1}^{n - 1} A_i} \quad \text{and} \quad \Prob\para{\bigcup_{i = 1}^{n - 1} \para{A_i \cap A_n}},
    \]
    and hence
    \[
        \Prob\para{\bigcup_{i = 1}^{n - 1} A_i} = \sum_{k = 1}^{n - 1} \para{-1}^{k + 1} \sum_{1 \leq i_1 < \cdots < i_k \leq n - 1} \Prob\para{A_{i_1} \cap A_{i_2} \cap \cdots \cap A_{i_k}},
    \]
    and if we let \(B_i = A_i \cap A_n\),
    \[
        \Prob\para{\bigcup_{i = 1}^{n - 1} B_i} = \sum_{k = 1}^{n - 1} \para{-1}^{k + 1} \sum_{1 \leq i_1 < \cdots < i_k \leq n - 1} \Prob\para{A_{i_1} \cap A_{i_2} \cap \cdots \cap A_{i_k} \cap A_n}.
    \]

    This gives exactly as desired.
\end{proof}

\begin{remark}
    Let \(\Omega\) be a finite set, and \(\para{\Omega, \mathcal{F}, \Prob}\) be a probability space, where \(\Prob\) is the uniform measure
    \[
        \Prob\para{A} = \frac{\abs{A}}{\abs{\Omega}}.
    \]

    Taking \(A_1, A_2, \ldots, A_n \in \mathcal{F}\), applying \autoref{thm:inclusion-exclusion} (Inclusion--Exclusion Formula) to \(\Prob\para{\bigcup_{i = 1}^{n} A_i}\), we have
    \[
        \abs{A_1 \cup A_2 \cup \cdots \cup A_n} = \sum_{k = 1}^{n} \para{-1}^{k + 1} \sum_{1 \leq i_1 < \cdots < i_n \leq n} \abs{A_{i_1} \cap \cdots \cap A_{i_n}},
    \]
    the Inclusion--Exclusion Principle on cardinalities. So the measure version implies the cardinality one.
\end{remark}

Now, if we truncate the above sums, in the small cases, we see for \(n = 2\)
\[
    \Prob\para{A \cup B} \leq \Prob\para{A} + \Prob\para{B},
\]
and for \(n = 3\),
\[
    \Prob\para{A \cup B \cup C} \geq \Prob\para{A} + \Prob\para{B} + \Prob\para{C} - \Prob\para{A \cap B} - \Prob\para{A \cap C} - \Prob\para{B \cap C}
\]
but
\[
    \Prob\para{A \cup B \cup C} \leq \Prob\para{A} + \Prob\para{B} + \Prob\para{C}.
\]

This leads to \autoref{prop:bonferroni-inequality}, the Bonferroni Inequalities.

\begin{proposition}[Bonferroni Inequalities]
    \label{prop:bonferroni-inequality}%
    We have
    \[
        \Prob\para{\bigcup_{i = 1}^{n} A_i} \begin{cases}
            \leq \sum_{k = 1}^{n} \para{-1}^{k + 1} \sum_{1 \leq i_1 < \cdots < i_k \leq n} \Prob\para{A_{i_1} \cap A_{i_2} \cap \cdots \cap A_{i_k}}, & r \text{ is odd}, \\
            \geq \sum_{k = 1}^{n} \para{-1}^{k + 1} \sum_{1 \leq i_1 < \cdots < i_k \leq n} \Prob\para{A_{i_1} \cap A_{i_2} \cap \cdots \cap A_{i_k}}, & r \text{ is even}.
        \end{cases}
    \]
\end{proposition}

\begin{proof}
    The base case is true for \(n = 2\). We assume these inequalities hold for \(n - 1\) events.

    If \(r\) is odd, let \(B_i = A_i \cap A_n\), and we have
    \[
        \Prob\para{\bigcup_{i = 1}^{n} A_i} = \Prob\para{\bigcup_{i = 1}^{n - 1} A_i} + \Prob\para{A_n} - \Prob\para{B_1 \cup B_2 \cup \cdots \cup B_{n - 1}}.
    \]

    We apply the induction hypothesis to \(r\) is odd for \(A_i\), giving
    \[
        \Prob\para{\bigcup_{i = 1}^{n - 1} A_i} \leq \sum_{k = 1}^{r} \para{-1}^{k + 1} \sum_{1 \leq i_1 < \cdots < i_n \leq n - 1} \Prob\para{A_{i_1} \cap \cdots \cap A_{i_n}},
    \]
    and to \(r - 1\) is even for \(B_i\), giving
    \[
        \Prob\para{\bigcup_{i = 1}^{n - 1} B_i} \leq \sum_{k = 1}^{r - 1} \para{-1}^{k + 1} \sum_{1 \leq i_1 < \cdots < i_n \leq n - 1} \Prob\para{A_{i_1} \cap \cdots \cap A_{i_n} \cap A_n}.
    \]

    Substituting these in finishes the proof for \(r\) is odd, and the proof is similar for \(r\) is even.
\end{proof}

\begin{example}[Counting Surjective Functions]
    \label{ex:counting-surjections}%
    Let
    \[
        \Omega = \brce{\functiondomain{f}{\brce{1, 2, \ldots, n}}{\brce{1, 2, \ldots, m}}},
    \]
    and take the subset
    \[
        A = \set{f \in \Omega}{f \text{ is surjective}}.
    \]

    We want to find \(\abs{A}\). We define
    \[
        A_i = \set{f \in \Omega}{i \notin \img f}.
    \]

    Then we have
    \[
        A = A_1\compl \cap A_2\compl \cap \cdots \cap A_m\compl = \para{A_1 \cup A_2 \cup \cdots \cup A_m}\compl.
    \]

    Now, we have
    \[
        \abs{A} = \abs{\Omega} - \abs{A_1 \cup \cdots \cup A_m},
    \]
    and by the Inclusion--Exclusion Principle, we have
    \begin{align*}
        \abs{A_1 \cup \cdots \cup A_m} &= \sum_{k = 1}^{m} \para{-1}^{k + 1} \sum_{1 \leq i_1 < \cdots < i_k \leq m} \abs{A_{i_1} \cap A_{i_2} \cap \cdots \cap A_{i_k}} \\
        &= \sum_{k = 1}^{m} \para{-1}^{k + 1} \binom{m}{k} \para{m - k}^n,
    \end{align*}
    so we have
    \[
        \abs{A} = \sum_{k = 0}^{m} \para{-1}^{k} \binom{m}{k} \para{m - k}^{n}.
    \]
\end{example}

\begin{example}[Derangements]
    \label{ex:derangements}%
    \emph{Derangements} are permutations with no fixed points. Let the event space be
    \[
        \Omega = \brce{\text{permutations of }\brce{1, 2, \ldots, n}}
    \]
    and the event of derangements be
    \[
        A = \brce{\text{derangements}} = \set{f \in \Omega}{\forall i = 1, 2, \ldots, n: f\para{i} \neq i}.
    \]

    What is the probability of \(A\)?

    For every \(i = 1, 2, \ldots, n\), we let
    \[
        A_i = \set{f \in \Omega}{f\para{i} = i}.
    \]

    Then, we have
    \[
        A = A_1\compl \cap A_2\compl \cap \cdots \cap A_n\compl = \para{\bigcup_{i = 1}^{n} A_i} \compl.
    \]

    Therefore, by the additivity of the probability measure, we have
    \[
        \Prob\para{A} = 1 - \Prob\para{\bigcup_{i = 1}^{n} A_i}.
    \]

    Using \autoref{thm:inclusion-exclusion} (Inclusion--Exclusion formula), we have
    \[
        \Prob\para{\bigcup_{i = 1}^{n} A_i} = \sum_{k = 1}^{n} \para{-1}^{k + 1} \sum_{1 \leq i_1 < \cdots < i_k \leq n} \Prob\para{A_{i_1} \cap A_{i_2} \cap \cdots \cap A_{i_k}}.
    \]

    We have
    \[
        \Prob\para{A_{i_1} \cap A_{i_2} \cap \cdots \cap A_{i_k}} = \frac{\para{n - k}!}{n!}
    \]
    since we have \(k\) points fixed in the permutations, and therefore \(\para{n - k}\) places free.
    
    Hence,
    \begin{align*}
        \Prob\para{\bigcup_{i = 1}^{n} A_i} &= \sum_{k = 1}^{n} \para{-1}^{k + 1} \binom{n}{k} \frac{\para{n - k}!}{n!} \\
        &= \sum_{k = 1}^{n} \para{-1}^{k + 1} \frac{1}{k!}.
    \end{align*}

    Hence, we have
    \[
        \Prob\para{A} = 1 - \sum_{k = 1}^{n} \para{-1}^{k + 1} \frac{1}{k!} = \sum_{k = 0}^{n} \frac{\para{-1}^{k}}{k!}.
    \]

    As \(n \to \infty\), we have the behaviour that
    \[
        \Prob\para{A} \to \frac{1}{e}.
    \]
\end{example}

\subsection{Independence and Conditional Probability}

Let \(\para{\Omega, \mathcal{F}, \Prob}\) be a probability space.

\begin{definition}[Independence]
    \label{def:independence}%
    We say two events \(A, B \in \mathcal{F}\) are \emph{independent}, denoted as \(A \indp B\), if
    \[
        \Prob\para{A \cap B} = \Prob\para{A} \cdot \Prob\para{B}.
    \]

    A \emph{countable collection} of events \(\para{A_n}_{n \in \NN}\) is said to be \emph{independent}, if for any (countable) collection of indices \(I = \brce{i_1, i_2, \ldots, i_k} \subseteq \NN\), we have
    \[
        \Prob\para{\bigcap_{i \in I} A_i} = \prod_{i \in I} \Prob\para{A_i}.
    \]
\end{definition}

\begin{remark}
    Pairwise independence \textbf{does not imply} independence.
\end{remark}

\begin{example}[Pairwise Independence does not imply Independence]
    \label{ex:pairwise-indep-does-not-imply-indep}
    We toss a fair coin twice. The sample space is
    \[
        \Omega = \brce{\para{0, 0}, \para{0, 1}, \para{1, 0}, \para{1, 1}}
    \]
    equipped with probability measure
    \[
        \Prob\para{\omega} = \frac{1}{4}
    \]  
    for all \(\omega \in \Omega\). Take
    \[
        A = \brce{\para{0, 0}, \para{0, 1}}, B = \brce{\para{0, 0}, \para{1, 0}}, C = \brce{\para{1, 0}, \para{0, 1}}.
    \]

    We have
    \[
        \Prob\para{A} = \Prob\para{B} = \Prob\para{C} = \frac{1}{2},
    \]
    and
    \[
        \Prob\para{A \cap B} = \Prob\para{A \cap C} = \Prob\para{B \cap C} = \frac{1}{4},
    \]
    so \(\brce{A, B, C}\) are \emph{pairwise} independent. But they are not \emph{independent}, since
    \[
        \Prob\para{A \cap B \cap C} = \Prob\para{\emptyset} = 0.
    \]
\end{example}

\begin{proposition}[Independence implies Independence of Complement]
    \label{prop:indep-compl}%
    If \(A \indp B\), then \(A \indp B\compl\).
\end{proposition}

\begin{proof}
    Investigating the probability, we have
    \begin{align*}
        \Prob\para{A \cap B\compl} &= \Prob\para{A} - \Prob\para{A \cap B} \\
        &= \Prob\para{A} - \Prob\para{A} \cdot \Prob\para{B} \\
        &= \Prob\para{A} \para{1 - \Prob\para{B}} \\
        &= \Prob\para{A} \Prob\para{B\compl}.
    \end{align*}
\end{proof}

An idea closely related is conditional probability.

\begin{definition}[Conditional Probability]
    \label{def:conditional-prob}%
    Let \(B \in \mathcal{F}\) with non-zero measure, \(\Prob\para{B} > 0\).

    For any \(A \in \mathcal{F}\), we define the conditional probability of \(A\) given \(B\), is given by
    \[
        \Prob\Conditional{A}{B} = \frac{\Prob\para{A \cap B}}{\Prob\para{B}}.
    \]
\end{definition}

If \(A \indp B\), then we have
\[
    \Prob\Conditional{A}{B} = \Prob\para{A}.
\]

\begin{proposition}[Countable Additivity Property for Conditional Probability]
    \label{prop:countable-additivity-conditional}%
    Let \(\para{A_n}_{n \in \NN}\) be a disjoint sequence of events in \(\mathcal{F}\). Then we have
    \[
        \Prob\Conditional{\bigcup_n A_n}{B} = \sum_{n} \Prob\Conditional{A_n}{B}.
    \]
\end{proposition}

\begin{proof}
    By definition, we have
    \begin{align*}
        \Prob\Conditional{\bigcup_{n} A_n}{B} &= \frac{\Prob\para{\bigcup_n A_n \cap B}}{\Prob\para{B}} \\
        &= \frac{\Prob\para{\bigcup_n \para{A_n \cap B}}}{\Prob\para{B}} \\
        &= \sum_{n} \frac{\Prob\para{A_n \cap B}}{\Prob{B}} \\
        &= \sum_{n} \Prob\Conditional{A_n}{B}.
    \end{align*}
\end{proof}

\begin{theorem}[Law of Total Probability]
    \label{thm:law-total-probability}%
    Suppose that the events \(\para{B_n}_{n \in \NN}\) is a \emph{partition} of \(\Omega\), i.e. pairwise disjoint with union \(\Omega\), with positive measure. i.e. \(\Prob\para{B_n} > 0\).

    Let \(A \in \mathcal{F}\). We have
    \[
        \Prob\para{A} = \sum_{n} \Prob\Conditional{A}{B_n} \Prob\para{B_n}.
    \]
\end{theorem}

\begin{proof}
    We have
    \begin{align*}
        \Prob\para{A} &= \Prob\para{A \cap \bigcup_{n} B_n} \\
        &= \Prob\para{\bigcup_n \para{B_n \cap A}} \\
        &= \sum_{n} \Prob\para{B_n \cap A} \\
        &= \sum_{n} \Prob\Conditional{A}{B_n} \Prob\para{B_n}.
    \end{align*}
\end{proof}

\begin{theorem}[Bayes' Formula]
    \label{thm:bayes}%
    Let \(\para{B_n}\) be a partition of \(\Omega\) with positive measure. We have
    \[
        \Prob\Conditional{B_n}{A} = \frac{\Prob\Conditional{A}{B_n} \cdot \Prob\para{B_n}}{\sum_{k} \Prob\Conditional{A}{B_k} \cdot \Prob\para{B_k}}.
    \]
\end{theorem}

\begin{proof}
    Indeed, we have
    \begin{align*}
        \Prob\Conditional{B_n}{A} &= \frac{\Prob\para{B_n \cap A}}{\Prob\para{A}} \\
        &= \frac{\Prob\Conditional{A}{B_n} \cdot \Prob\para{B_n}}{\Prob\para{A}},
    \end{align*}
    and we apply \autoref{thm:law-total-probability} (Law of Total Probability) for the numerator.
\end{proof}

This formula is the basis of Bayesian statistics. If we know the probability of \(\Prob\para{B_n}\), and we have a model which gives us \(\Prob\Conditional{A}{B_k}\), with this formula we can find the posterior probability of \(B_n\) given that \(A\) occurs.

\begin{example}[False Positives for a Rare Condition]
    \label{ex:false-positives}%
    Suppose that a condition A affects \(0.1\%\) of the population. We have a medical test which is positive for \(98\%\) of the affected population, and \(1\%\) of those unaffected.

    Pick an individual at random. What is the probability they suffer from A given they tested positive?

    We define the events
    \[
        A = \brce{\text{individuals suffering from A}},
    \]
    and
    \[
        P = \brce{\text{individuals tests positive}}.
    \]

    We want to find \(\Prob\Conditional{A}{P}\). The given conditions state that
    \[
        \Prob\para{A} = 0.001, \quad \Prob\Conditional{P}{A} = 0.98, \quad \text{and} \quad \Prob\Conditional{P}{A\compl} = 0.01.
    \]

    Hence, using \autoref{thm:bayes} (Bayes' Formula), we have
    \begin{align*}
        \Prob\Conditional{A}{P} &= \frac{\Prob\Conditional{P}{A} \cdot \Prob\para{A}}{\Prob\Conditional{P}{A} \cdot \Prob\para{A} + \Prob\Conditional{P}{A\compl} \cdot \Prob\para{A\compl}} = 0.089 \approx 0.09
    \end{align*}
    so
    \[
        \Prob\Conditional{A}{P} \approx 0.09.
    \]

    This is counterintuitive. We have that
    \[
        \Prob\Conditional{A}{P} = \frac{1}{1 + \frac{\Prob\Conditional{P}{A\compl} \cdot \Prob\para{A\compl}}{\Prob\Conditional{P}{A} \cdot \Prob\para{A}}}.
    \]

    We have \(\Prob\para{A\compl}\) and \(\Prob\Conditional{P}{A}\) are both very close to \(1\), so the relevant ratio simplifies to \(\frac{\Prob\Conditional{P}{A\compl}}{\Prob\para{A}}\). But
    \[
        \Prob\Conditional{P}{A\compl} \gg \Prob\para{A}
    \]
    and this makes \(\Prob\Conditional{A}{P}\) is \emph{very small}.

    Let's put in some concrete numbers. Suppose we have a population of \(1000\) people, with \(1\) suffering from A. Among the \(999\) not suffering, about \(10\) people will test positive. In total, there will be \(11\) people testing positive. Among those we pick one at random. So the probability of being positive, given they actually test positive, is only \(\frac{1}{11}\).
\end{example}

\begin{example}[Children and Thursday]
    \label{ex:children-thursday}%
    \begin{enumerate}
        \item What is the probability that I have \(2\) boys, given I have \(2\) children, one of whom is a boy?
    
            Since no further information is given, we take all outcomes to be equally likely.

            We denote the events
            \[
                \mathrm{BG} = \brce{\text{elder child is a boy, younger is a girl}},
            \]
            and the events \(\mathrm{GB}, \mathrm{GG}, \mathrm{BB}\). We want to find
            \[
                \Prob\Conditional{\mathrm{BB}}{\mathrm{BB} \cup \mathrm{BG} \cup \mathrm{GB}} = \frac{\Prob\para{\mathrm{BB}}}{\Prob\para{\mathrm{BB} \cup \mathrm{BG} \cup \mathrm{GB}}} = \frac{\frac{1}{4}}{\frac{3}{4}} = \frac{1}{3}.
            \]

        \item What is the probability that I have \(2\) boys, given the eldest one is a boy?
        
            We have
            \[
                \Prob\Conditional{\mathrm{BB}}{\mathrm{BB} \cup \mathrm{BG}} = \frac{1}{2}.
            \]

        \item What is the probability that I have \(2\) boys, given that one of them is a boy born on a Thursday?
        
        Let
        \[
            \mathrm{TG} = \brce{\text{elder boy is a boy on Thursday, younger is a girl}},
        \]
        and \(\mathrm{GT}\) similarly, and
        \[
            \mathrm{TN} = \brce{\text{elder boy is a boy on Thursday, younger is a boy not born on Thursday}},
        \]
        and \(\mathrm{NT, TT}\) similarly. We want to find
        \begin{align*}
            &\phantom{=} \Prob\Conditional{\mathrm{TT} \cup \mathrm{TN} \cup \mathrm{NT}}{\mathrm{TT} \cup \mathrm{TG} \cup \mathrm{GT} \cup \mathrm{TN} \cup \mathrm{NT}} \\
            &= \frac{\Prob\para{\mathrm{TT} \cup \mathrm{TN} \cup \mathrm{NT}}}{\Prob\para{\mathrm{TT} \cup \mathrm{TG} \cup \mathrm{GT} \cup \mathrm{TN} \cup \mathrm{NT}}} \\
            &= \frac{\frac{1}{2} \cdot \frac{1}{7} \cdot \frac{1}{2} \cdot \frac{1}{7} + \frac{1}{2} \cdot \frac{1}{7} \cdot \frac{1}{2} \cdot \mathbf{\frac{6}{7}} + \frac{1}{2} \cdot \mathbf{\frac{6}{7}} \cdot \frac{1}{2} \cdot \frac{1}{7}}{\frac{1}{2} \cdot \frac{1}{7} \cdot \frac{1}{2} \cdot \frac{1}{7} + \frac{1}{2} \cdot \frac{1}{7} \cdot \frac{1}{2} + \frac{1}{2} \cdot \frac{1}{7} \cdot \frac{1}{2} + \frac{1}{2} \cdot \frac{1}{7} \cdot \frac{1}{2} \cdot \mathbf{\frac{6}{7}} \cdot 2}
 \\
            &= \frac{13}{27}.
        \end{align*}
    \end{enumerate}
\end{example}

\begin{table}[ht!]
    \centering
    \begin{tabular}{cccc}
        All Applicants & Admitted & Rejected & \% Admitted \\
        \hline
        State & 25 & 25 & 50\% \\
        Independent & 28 & 22 & 56\%
    \end{tabular}
    
    \vspace{1em}

    \begin{tabular}{cccc}
        London Schools & Admitted & Rejected & \% Admitted \\
        \hline
        State & 15 & 22 & 41\% \\
        Independent & 5 & 8 & 38\%
    \end{tabular}
    
    \vspace{1em}

    \begin{tabular}{cccc}
        Cambridge Schools & Admitted & Rejected & \% Admitted \\
        \hline
        State & 10 & 3 & 77\% \\
        Independent & 23 & 14 & 62\%
    \end{tabular}

    \caption{Data of Admissions}
    \label{tab:simpson-paradox-data}
\end{table}

\begin{example}[Simpson's Paradox]
    \label{ex:simpson-paradox}%
    We refer to \autoref{tab:simpson-paradox-data}. This phenomenon is called \emph{cofounding} in statistics. It arises when we aggregate data from disparate populations.

    Let
    \begin{align*}
        A &= \brce{\text{individual is admitted}}, \\
        B &= \brce{\text{individual is from London}}, \\
        B\compl &= \brce{\text{individual is from Cambridge}}, \\
        C &= \brce{\text{state school}}, \\
        C\compl &= \brce{\text{independent school}}.
    \end{align*}

    We notice that from the latter two tables, that
    \begin{align*}
        \Prob\Conditional{A}{B \cap C} &> \Prob\Conditional{A}{B \cap C \compl}, \\
        \Prob\Conditional{A}{B\compl \cap C} &> \Prob\Conditional{A}{B\compl \cap C \compl}.
    \end{align*}

    However, we have from the first table, that
    \[
        \Prob\Conditional{A}{C} < \Prob\Conditional{A}{C\compl}.
    \]

    Observe that from countable additivity of conditional probability,
    \begin{align*}
        \Prob\Conditional{A}{C} &= \Prob\Conditional{A \cap B}{C} + \Prob\Conditional{A \cap B\compl}{C} \\
        &= \frac{\Prob\para{A \cap B \cap C}}{\Prob\para{C}} + \frac{\Prob\para{A \cap B\compl \cap C}}{\Prob\para{C}} \\
        &= \frac{\Prob\Conditional{A}{B \cap C} \cdot \Prob{B \cap C}}{\Prob\para{C}} + \frac{\Prob\Conditional{A}{B \compl \cap C} \cdot \Prob\para{B\compl \cap C}}{\Prob\para{C}} \\
        &= \Prob\Conditional{A}{B \cap C} \cdot \Prob\Conditional{B}{C} + \Prob\Conditional{A}{B\compl \cap C} \cdot \Prob\Conditional{B \compl}{C} \\
        &> \Prob\Conditional{A}{B \cap C\compl} \cdot \Prob\Conditional{B}{C} + \Prob\Conditional{A}{B\compl \cap C\compl} \cdot \Prob \Conditional{B\compl}{C}.
    \end{align*}

    Now, note that if \(\Prob\Conditional{B}{C} = \Prob\Conditional{B}{C\compl}\), which is not valid here, then we would get
    \begin{align*}
        \Prob\Conditional{A}{C} &> \Prob\Conditional{A}{B \cap C\compl} \cdot \Prob\Conditional{B}{C} + \Prob\Conditional{A}{B\compl \cap C\compl} \cdot \Prob \Conditional{B\compl}{C} \\
        &= \Prob\Conditional{A}{B \cap C\compl} \cdot \Prob\Conditional{B}{C\compl} + \Prob\Conditional{A}{B\compl \cap C\compl} \cdot \Prob \Conditional{B\compl}{C\compl} \\
        &= \Prob\Conditional{A}{C\compl}
    \end{align*}
    hence this does not hold.

    Effectively, we are dealing with inequalities
    \[
        \frac{a_1}{a_2} > \frac{b_1}{b_2} \quad \text{and} \quad \frac{c_1}{c_2} > \frac{d_1}{d_2},
    \]
    but we cannot say anything about
    \[
        \frac{a_1 + c_1}{a_2 + c_2} \quad \text{and} \quad \frac{b_1 + d_1}{b_2 + d_2}.
    \]
\end{example}